\section{What the model is for}

GGM was developed by Ruud Egging, building on the European Gas Model and the World Gas Model \citep[p.~8]{egging2019global}, and has since been applied and documented together with Franziska Holz at DIW Berlin.
It belongs to a family of models that represent natural gas markets as a game between strategic players and compute the resulting market equilibrium \citep{gabriel2005mixed, egging2008complementarity, egging2010world}.
Such equilibrium models followed a long period of optimisation models of energy markets, because they can represent a market game between players rather than a single planner minimising cost; applied to European gas, they confirmed an oligopolistic market structure in which a few large exporters can influence prices \citep{holz2015modeling, holz2008strategic}.

That is why GGM represents market power.
Gas markets are shaped by a small number of exporting countries and by long-lived infrastructure, and the questions the model is used for concern both.
It has been applied to the future of global gas under climate constraints \citep{holz2015global}, to infrastructure projects of common European interest and to security of supply \citep[p.~8]{egging2019global}, and most recently to whether the European Union still needs Russian gas under different climate pathways \citep{barner2025russian, holz2024diw}.
Its results are long-term scenarios rather than forecasts: they show how production, trade, infrastructure and prices would develop if the assumed demand, supply and policies held.
The model covers gas alone, without long-term contracts or oil-price indexation, so competition with other fuels enters only through the shape of the demand curves; results have to be read with these limits in mind.
